%%%%

%%%   Самарский государственный технический университет

%%%   Кафедра прикладной математики и информатики

%%%

%%%   Преамбула для статьи в журнал "Вестник Самарского государственного

%%%   технического университета. Серия: «Физико-математические науки»"

%%%   Ориентирована под MikTEx 2.4 for Win32

%%%

%%%   Хотя сам журнал собирается под GNU/Linux на дистрибутиве TeXLive



\documentclass[11pt,a4paper,twoside]{article}



\usepackage[cp1251]{inputenc}

\usepackage[T2A]{fontenc}

\usepackage[english,russian]{babel}

\usepackage{samgtubib}

\usepackage[dvips]{graphicx}

\usepackage[multiple]{footmisc}



\usepackage{samgtu}

\renewcommand{\VestnikYear}{2009} % Год издания Вестника

\renewcommand{\VestnikNumber}{2\,(19)} % Номер Вестника

\renewcommand{\VestnikMonth}{Ноябрь} % Месяц выхода Вестника

\renewcommand{\VestnikData}{01.11.09} % Дата подписания Вестника в печать

\renewcommand{\VestnikUPL}{26,64} % Усл. печ. л.

\renewcommand{\VestnikUIL}{25,73} % Уч.-изд. л.

\renewcommand{\VestnikRegNumber}{479/09} % Регистрационный номер

\renewcommand{\VestnikOrder}{994} % Номер заказа







%

%\begin{document}



\renewcommand{\firstpage}{160}

\renewcommand{\lastpage}{164}







\udk{517.968.4}

\msc{45G10}



\title{Нелинейные уравнения с~весовыми операторами типа потенциала в~пространствах Лебега}

{Нелинейные уравнения с весовыми операторами типа потенциала \dots}



\author{С.\,Н.~Асхабов}{А\,с\,х\,а\,б\,о\,в~С.\,Н.}



\institute{Чеченский государственный университет, \\

364907, Грозный, ул. А.~Шерипова, 32.}





\email{E-mail}{askhabov@yandex.ru}



\otherinf{\fullauthor{Султан Нажмудинович Асхабов}\regalia{д.ф.-м.н., доц.}\position{декан}

\dept{факультет математики и компьютерных технологий}}



\keywords{нелинейные уравнения, оператор типа потенциала, монотонный оператор}



\workabstract{Методом монотонных операторов для различных классов

нелинейных уравнений с~весовыми операторами типа потенциала

доказаны глобальные теоремы о~существовании, единственности и~оценках решений в~пространствах Лебега.}



\engtitle{Nonlinear equations with weighted potential type operators in Lebesgue spaces}



\engauthor{S.\,N.~Askhabov}{A\,s\,k\,h\,a\,b\,o\,v~S.\,N.}





\enginstitute{Chechen State University,\\

32, A.~Sheripova st., Groznyi, 364907, Russia.}





\otherenginf{\fullauthor{Sultan N. Askhabov} \regalia{Dr.\,Sci.

(Phys. \& Math.)} \position{Dean of Faculty} \dept{Faculty of Mathematics and Computer Technology} }



\engabstract{By method of monotone operators, existence and

uniqueness theorems are proved for some classes of nonlinear

equations with weighted potential type operators in Lebesgue

spaces.}



\engkeywords{nonlinear equations, potential type operator, monotone operator}



\date{28/VI/2011}

\revisiondate{27/VII/2011}



\maketitle



\small



В вещественных пространствах $L_p(R^1)=L_p(-\infty, \infty)$, $1<p<\infty$, рассматриваются нелинейные интегральные

уравнения вида

\begin{gather}

\label{equat1}

u(x)+\int _{-\infty}^{\infty}\frac{[a(x)-a(t)]\,F(t,

u(t))}{|x-t|^{1-\alpha}}dt=f(x),\\

 \label{equat2} u(x)+F\left(x,

\int _{-\infty}^{\infty}\frac{[a(x)-a(t)]\,u(t)}{|x-t|^{1-\alpha}}dt\right)=f(x),

\end{gather}

для которых методом монотонных (по Браудеру"--~Минти) операторов

(см., например,~\cite{askhabov:ASK}) доказываются глобальные

теоремы о существовании, единственности и оценках решений.



Для упрощения записей введём следующие обозначения:

\begin{gather*}

L_p(R^1)=L_p, \quad \|\cdot \|_{L_p(R^1)}=\|\cdot \|_p, \quad

p'=\frac{p}{p-1}, \quad \langle u,

v\rangle=\int _{-\infty}^{\infty}u(x)\,v(x)\,dx,\\

(I^{\alpha}u

)(x)=\int _{-\infty}^{\infty}\frac{u(t)\,dt}{|x-t|^{1-\alpha}},

\qquad

(A^{\alpha}u)(x)=\int _{-\infty}^{\infty}\frac{[a(x)-a(t)]\,u(t)}{|x-t|^{1-\alpha}}dt.

\end{gather*}



В силу известной теоремы Харди"--~Литтлвуда (см.,~например,~\cite{askhabov:ASK}) оператор типа потенциала

$I^\alpha$ действует непрерывно из $L_p$ в $L_{p/(1-\alpha\,p)}$,

если $0<\alpha<1$ и $1<p<1/\alpha$, причём

\begin{equation}

\label{hardi1} \|I^\alpha u\|_{p/(1-\alpha\,p)}\leq

\|I^\alpha\|_{p\to p/(1-\alpha\,p)}\|u\|_p \qquad \forall u\in

L_p,

\end{equation}

где $\|I^\alpha\|_{p\to p/(1-\alpha\,p)}$ есть норма оператора $I^\alpha : L_p\to L_{p/(1-\alpha\,p)}$.



Справедливы следующие (двойственные) леммы.



\begin{lemma}[1]Пусть $0<\alpha<1,$ $2/(1+\alpha)<p<1/\alpha$ и $a\in

L_{p/[p(1+\alpha)-2]}.$ Тогда оператор $A^{\alpha}$ действует

непрерывно из $L_p$ в $L_{p'},$ причём

\begin{gather}

\label{askh1} \|A^{\alpha}u\|_{p'}\leq  2\,\|I^\alpha\|_{p\to

p/(1-\alpha\,p)}\|a\|_{p/[p(1+\alpha)-2]}\|u\|_p ,\\

\label{posit1} \langle A^\alpha u, u\rangle = 0 \qquad \forall

u(x)\in L_p.

\end{gather}

\end{lemma}



\begin{lemma}[2]Пусть $0<\alpha<1,$ $1/(1-\alpha)<p<2/(1-\alpha)$ и $a\in

L_{p/[2-p(1-\alpha)]}.$ Тогда оператор $A^{\alpha}$ действует

непрерывно из $L_{p'}$ в $L_p,$ причём

\begin{gather}

\label{askh2} \|A^{\alpha}u\|_p\leq

2\,\|I^\alpha\|_{p/(1+\alpha\,p)\to

p}\|a\|_{p/[2-p(1-\alpha)]}\|u\|_{p'},\\

\label{posit2} \langle A^\alpha u, u\rangle = 0 \qquad \forall

u(x)\in L_{p'}.

\end{gather}

\end{lemma}



Лемма 1 доказана в [\citen{askhabov:3}]. Докажем  лемму~2.



\begin{proof}

Пусть $u\in L_{p'}$. Тогда, применяя неравенство Гёльдера с показателями $(1+\alpha\,p)/(p-1)$ и $(1+\alpha\,p)/[2-p(1-\alpha)]$, имеем

\begin{equation}

\label{multipl} \|a\cdot u\|_{p/(1+\alpha\,p)}\leq

\|a\|_{p/[2-p(1-\alpha)]}\|u\|_{p'}.

\end{equation}

Итак, $a\cdot u\in L_{p/(1+\alpha\,p)}$. 

Так как $1<p/(1+\alpha\,p)<1/\alpha$ (первое неравенство равносильно условию, что $1/(1-\alpha)<p$, а второе "--- очевидно), 

согласно теореме Харди"--~Литтлвуда $I^\alpha(a\,u)\in L_p$, поскольку

$${\frac{\frac{p}{1+\alpha\,p}}{1-\frac{\alpha\,p}{1+\alpha\,p}}=p},$$

причём $$\|I^\alpha(a\cdot u)\|_p\leq

\|I^\alpha\|_{p/(1+\alpha\,p)\to p}\|a\cdot

u\|_{p/(1+\alpha\,p)}.$$ 

Воспользовавшись оценкой \eqref{multipl},

из последнего неравенства получаем

\begin{equation}

\label{hardi2}

\|I^\alpha(a\cdot u)\|_p\leq \|I^\alpha\|_{p/(1+\alpha\,p)\to p}\|a\|_{p/[2-p(1-\alpha)]}\|u\|_{p'} .

\end{equation}



Обозначим $I^\alpha u=v$. Так как $u\in L_{p'}$ и $1<p'<1/\alpha$ (первое неравенство очевидно, а второе равносильно условию, что

$1/(1-\alpha)<p$), cогласно теореме Харди"--~Литтлвуда

$I^\alpha u\in L_q$, 

где

$${q=\frac{p'}{1-\alpha\,p'}=\frac{p}{p(1-\alpha)-1}},$$

т.\,е. $v=I^\alpha u\in L_{p/[p(1-\alpha)-1]}$, $\forall u\in

L_{p'}$, причём, в силу неравенства \eqref{hardi1},

\begin{equation}

\label{hardi3}

\|I^\alpha u\|_{p/[p(1-\alpha)-1]}\leq \|I^\alpha\|_{p/(p-1)\to p/[p(1-\alpha)-1]}\|u\|_{p'} .

\end{equation}

Далее, применяя неравенство Гёльдера с показателями

$1/[p(1{-}\alpha){-}1]$ и \mbox{$1/[2{-}p(1{-}\alpha)]$}, имеем

$$

\|a\cdot

I^\alpha

u\|_p=\left(\int _{-\infty}^{\infty}|a(x)|^p|v(x)|^pdx\right)^{1/p}\leq

\|a\|_{p/[2-p(1-\alpha)]}\|I^\alpha u\|_{p/[p(1-\alpha)-1]}.

$$

Поэтому с учётом оценки \eqref{hardi3} сразу получаем

\begin{equation}

\label{hardi4}

\|a\cdot I^\alpha u\|_p\leq \|I^\alpha\|_{p/(p-1)\to p/[p(1-\alpha)-1]}\|a\|_{p/[2-p(1-\alpha)]}\|u\|_{p'} .

\end{equation}

Так как $A^\alpha u=a\cdot I^\alpha u-I^\alpha(a\cdot u)$, из неравенств  \eqref{hardi2} и \eqref{hardi4} следует, что оператор $A^\alpha$ действует непрерывно из $L_{p'}$ в $L_p$. 

Используя для оценки $\|A^\alpha u\|_p$ сначала неравенство Минковского, затем оценки \eqref{hardi2}, \eqref{hardi4} и~очевидное (см.,~например, \cite[c.~247]{askhabov:KOLM}) равенство $\|I^\alpha\|_{p\to

p/(1-\alpha\,p)}=\|I^\alpha\|_{p/[p(1+\alpha)-1]\to p'}$

(поскольку $I^\alpha$ самосопряжённый оператор), легко получаем неравенство \eqref{askh2}).



Осталось доказать равенство \eqref{posit2}. 

Так как оператор $I^\alpha$ является симметрическим

$$

\langle A^\alpha u, u\rangle=\langle a\,I^\alpha u,

u\rangle-\langle I^\alpha (a\,u), u\rangle=\langle I^\alpha u,

a\,u\rangle-\langle a\,u, I^\alpha u\rangle=0,

$$

что и требовалось доказать.

\end{proof}



\smallskip



Приступим теперь к исследованию нелинейного уравнения

(\ref{equat1}), содержащего оператор $A^\alpha$. Обозначим через

$L_p^+$ множество всех неотрицательных функций из $L_p$. Всюду

далее предполагается, что функция $F(x,\,t)$, порождающая оператор

Немыцкого $Fu=F[x,\,u(x)]$, определена при $x\in \mathbb R^1$, $t\in \mathbb R^1$

и удовлетворяет условиям Каратеодори: она измерима по $x$ при

каждом фиксированном $t$ и непрерывна по $t$ почти для всех $x$.





\begin{theorem}[1]Пусть $0<\alpha<1,$ $1/(1-\alpha)<p<2/(1-\alpha)$ и $a\in

L_{p/[2-p(1-\alpha)]}.$

Если нелинейность $F(x,\,t)$ удовлетворяет следующим условиям{\/\rm:}



\begin{itemize}

\item[{\rm 1)}] $|F(x,\,t)|\leq c(x)+d_1\,|t|^{p-1},$ где $c(x)\in L_{p'}^+,$ $d_1>0;$



\item[{\rm 2)}] $F(x,\,t)$ строго возрастает по $t$ почти при каждом фиксированном $x;$



\item[{\rm 3)}] $F(x,\,t)\cdot t\geq d_2 |t|^p-D(x),$ где $D(x)\in L_1^+,$ $d_2>0,$

\end{itemize}

то уравнение \eqref{equat1} имеет единственное решение

$u^*\in L_p$ при любом $f\in L_p.$

Кроме того$,$ если условия $1)$ и $3)$ выполнены при $c(x)=D(x)=0,$ то справедлива оценка

\begin{equation}

\label{estim1}

\|u^*\|_p\leq d_1d_2^{-1}\,\|f\|_p.

\end{equation}

\end{theorem}



\begin{proof}

Из условий 1)--3) в силу теорем 2.1, 2.2 и оценки (2.3) из \cite{askhabov:ASK} вытекает, что оператор Немыцкого $F$, порождённый функцией $F(x,\, t)$, отображает пространство $L_p$ на сопряжённое с~ним пространство $L_{p'}$, непрерывен, строго

монотонен и коэрцитивен. 

Поэтому в силу леммы 2.1 из \cite{askhabov:ASK} существует обратный оператор $F^{-1}$, отображающий $L_{p'}$ на $L_p$, хеминепрерывный, строго монотонный и коэрцитивный. 

Из леммы~2 вытекает, что оператор $A^\alpha$ действует из $L_{p'}$ в $L_p$, непрерывен и положителен. 

Значит, оператор $\Phi=F^{-1}+A^\alpha$ отображает $L_{p'}$ на $L_p$, хеминепрерывен, строго монотонен и коэрцитивен. Следовательно, по теореме Браудера"--~Минти (см. теорема~1.1 из \cite{askhabov:ASK}), уравнение $F^{-1}v+A^\alpha v=f$ имеет единственное решение $v^*\in L_{p'}$. 

Но тогда непосредственно проверяется, что $u^*=F^{-1}v^*\in L_p$ является решением уравнения $u+A^\alpha Fu=f$, т.\,е. данного уравнения (\ref{equat1}), и это решение $u^*$ является единственным в $L_p$ (что легко доказывается методом от противного).



Осталось доказать оценку \eqref{estim1}. Используя условия 1) и 3)

при $c(x)=D(x)=0$, равенство \eqref{posit2} и равенство

$u^*+A^\alpha Fu^*=f$, имеем

\begin{multline*}

d_2\|u^*\|_p^p\leq \langle u^*, Fu^*\rangle=\langle  u^*+A^\alpha

Fu^*, Fu^*\rangle =\langle f, Fu^*\rangle \leq \\ \leq d_1\|f\|_p\|Fu^*\|_{p'}\leq

d_1\|f\|_p\|u^*\|_p^{p-1},

\end{multline*}

 откуда непосредственно получаем оценку

(\ref{estim1}).

\end{proof}



\smallskip



\begin{newthm}{Следствие 1}Если $1/2<\alpha<3/4$ и $a\in L_{2/(2\alpha-1)},$ то при любом

$f\in L_4$ уравнение

$$

u(x)+\int _{-\infty}^\infty

\frac{[a(x)-a(t)]\,u^3(t)}{|x-t|^{1-\alpha}}dt=f(x)

$$

имеет единственное решение $u^*\in L_4,$ причём $\|u^*\|_4\leq \|f\|_4.$

\end{newthm}



\smallskip



\begin{newthm}{Следствие 2}Если $0<\alpha<1/4$ и $a\in L_{2/(1+2\alpha)},$ то при любом $f\in L_{4/3}$ уравнение

$$

u(x)+\int _{-\infty}^\infty

\frac{[a(x)-a(t)]\,u^{1/3}(t)}{|x-t|^{1-\alpha}}dt=f(x)

$$

имеет единственное решение $u^*\in L_{4/3},$ причём $\|u^*\|_{4/3}\leq \|f\|_{4/3}.$

\end{newthm}



\smallskip



В следующей теореме, относящейся к нелинейному уравнению \eqref{equat2}, условия на нелинейность $F(x,\,t)$ подбираются

так, чтобы порождаемый ею оператор Немыцкого $F$ действовал непрерывно из $L_{p'}$ в $L_p$ и был строго монотонным и

коэрцитивным.



\smallskip



\begin{theorem}[2]Пусть $0<\alpha<1,$ $2/(1+\alpha)<p<1/\alpha$ и $a\in L_{p/[p(1+\alpha)-2]}.$

Если нелинейность $F(x,\,t)$ удовлетворяет следующим условиям{\/\rm:}

\begin{itemize}



\item[{\rm 4)}] $|F(x,\,t)|\leq g(x)+d_3\,|t|^{1/(p-1)},$ где $g(x)\in L_p^+,$ $d_3>0;$



\item[{\rm 5)}] $F(x,\,t)$ строго возрастает по $t$ почти при каждом фиксированном $x;$



\item[{\rm 6)}] $F(x,\,t)\cdot t\geq d_4 |t|^{p/(p-1)}-D(x),$ где $D(x)\in L_1^+,$ $d_4>0,$



\end{itemize}

то уравнение \eqref{equat2} имеет единственное решение $u^*\in L_p$ при любом $f\in L_p$. Кроме того, если условия $4)$ и

$6)$ выполнены при $g(x)=D(x)=0,$ то справедлива оценка

\begin{equation}

\label{estim2}

\|u^*-f\|_p\leq \left(2\,d_3^pd_4^{-1}\|I^\alpha\|_{p\to p/(1-\alpha p)}\|a\|_{p/[p(1+\alpha)-2]}\|f\|_p\right)^{1/(p-1)}.

\end{equation}

\end{theorem}



\begin{proof}

Из леммы 1 и условий 4)--6) вытекает, соответственно, что весовой оператор типа потенциала $A^\alpha$ действует из $L_p$ в~$L_{p'}$, непрерывен и положителен, а оператор Немыцкого $F$ действует, наоборот, из $L_{p'}$ в $L_p$, непрерывен, строго

монотонен и коэрцитивен. 

Значит, по лемме~2.1 из \cite{askhabov:ASK}, существует хеминепрерывный, строго монтонный, коэрцитивный обратный оператор $F^{-1}$, действующий из $L_p$ в $L_{p'}$.



Запишем данное уравнение \eqref{equat2}  в операторном виде: $u+FA^\alpha u=f$. 

Полагая в нём $u=f-v$ и применяя затем к обеим частям получившегося уравнения оператор $F^{-1}$, приходим к

уравнению $\Phi v=A^\alpha f$, где $\Phi v=F^{-1}v+A^\alpha v$. 

В~силу указанных выше свойств операторов $A^\alpha$ и $F^{-1}$ оператор $\Phi$ действует из $L_p$ в $L_{p'}$, хеминепрерывен,

строго монотонен и коэрцитивен. 

Следовательно, по теореме Браудера"--~Минти (см.~теорему 1.1 из~\cite{askhabov:ASK}), уравнение $\Phi v=A^\alpha f$ имеет единственное решение $v^*\in L_p$. 

Но тогда непосредственно проверяется, что уравнение $u+FA^\alpha u=f$, т.\,е. данное уравнение \eqref{equat2}, имеет

решение $u^*=f-v^*\in L_p$. 

Единственность этого решения $u^*$ легко устанавливается методом от противного.



Осталось доказать оценку \eqref{estim2}. 

Воспользуемся равенствами $u^*+FA^\alpha u^*=f$, $F^{-1}v^*+A^\alpha v^*=A^\alpha f$, где $u^*=f-v^*$. 

Положим $w=F^{-1}v^*$. Тогда $Fw=v^*$. Используя условия 4) и 6) при $g(x)=D(x)=0$, неравенство \eqref{askh1}, равенство \eqref{posit1} и неравенство Гёльдера, имеем

\begin{multline*}

d_4\|w\|_{p'}^{p'}\leq \langle Fw, w\rangle=\langle v^*,

F^{-1}v^*\rangle=\langle v^*, F^{-1}v^*\rangle+\langle v^*,

A^{\alpha}v^*\rangle=\\

=\langle v^*, A^\alpha f\rangle\leq \|v^*\|_p\|A^\alpha

f\|_{p'}\leq 2\,\|I^\alpha\|_{p\to

p/(1-\alpha\,p)}\|a\|_{p/[p(1+\alpha)-2]}\|f\|_p \|v^*\|_p=\\

=2\,\|I^\alpha\|_{p\to

p/(1-\alpha\,p)}\|a\|_{p/[p(1+\alpha)-2]}\|f\|_p \|Fw\|_p\leq\\

\leq 2\,d_3\|I^\alpha\|_{p\to

p/(1-\alpha\,p)}\|a\|_{p/[p(1+\alpha)-2]}\|f\|_p

\|w\|_{p'}^{p'-1},

\end{multline*}

 откуда

\begin{equation}

\label{estim3}

\|w\|_{p'}\leq 2d_3d_4^{-1}\|I^\alpha\|_{p\to p/(1-\alpha\,p)}\|a\|_{p/[p(1+\alpha)-2]}\|f\|_p.

\end{equation}

Так как $\|f-u^*\|_p=\|v^*\|_p=\|Fw\|_p\leq \|w\|_{p'}^{p'-1}$, используя оценку \eqref{estim3} с учётом, что $p'-1=1/(p-1)$,

получаем

$$

\|u^*-f\|_p\leq d_3\left(2\,d_3d_4^{-1}\|I^\alpha\|_{p\to

p/(1-\alpha

p)}\|a\|_{p/[p(1+\alpha)-2]}\|f\|_p\right)^{1/(p-1)},

$$

что равносильно доказываемой оценке \eqref{estim2}.

\end{proof}



\smallskip



\begin{newthm}{Следствие 3}Если $0<\alpha<1/4$ и $a\in L_{2/(1+2\alpha)},$ то при любом $f\in

L_4$ уравнение

$$

u(x)+\left(\int _{-\infty}^{\infty}\frac{[a(x)-a(t)]\,u(t)}{|x-t|^{1-\alpha}}dt\right)^{1/3}=f(x)

$$

имеет единственное решение $u^*\in L_4,$ причём

$$

\|u^*-f\|_4\leq \left(2\,\|I^\alpha\|_{4\to

4/(1-4\alpha)}\|a\|_{2/(1+2\alpha)}\|f\|_4\right)^{1/3}.

$$

\end{newthm}



\smallskip



\begin{newthm}{Следствие 4}Если $1/2<\alpha<3/4$ и $a\in L_{2/(2\alpha-1)},$ то при любом

$f\in L_{4/3}$ уравнение

$$

u(x)+\left(\int _{-\infty}^{\infty}\frac{[a(x)-a(t)]\,u(t)}{|x-t|^{1-\alpha}}dt\right)^3=f(x)

$$

имеет единственное решение $u^*\in L_{4/3},$ причём

$$

\|u^*-f\|_{4/3}\leq \left(2\,\|I^\alpha\|_{4/3\to

4/(3-4\alpha)}\|a\|_{2/(2\alpha-1)}\|f\|_{4/3}\right)^3.

$$

\end{newthm}



Вопрос о приближённом решении уравнений (1) и (2) рассмотрен в [\citen{askhabov:4}].



\begin{thebibliography}{9}



\RBibitem{askhabov:ASK} 

\by Асхабов~С.\,Н. 

\book Нелинейные уравнения типа свёртки 

\publaddr М. 

\publ Физматлит 

\yr 2009

\totalpages 304

\misctransl

\by Askhabov~S.\,N.

\book Nonlinear equations of convolution type

\publaddr Moscow

\publ Fizmatlit

\yr 2009

\totalpages 304









\RBibitem{askhabov:KOLM} 

\by Колмогоров~А.\,Н., Фомин~С.\,В. 

\book Элементы теории функций и функционального анализа 

\publaddr М.

\publ Физматлит 

\yr 2004 

\totalpages 570

\misctransl

\book Elements of the Theory of Functions and Functional Analysis 

\by Kolmogorov~A.\,N.,  Fomin~S.\,V.

\publaddr Moscow

\publ Fizmatlit

\yr 2004 

\totalpages 570



\RBibitem{askhabov:3} 

\by Асхабов~С.\,Н.

\paper Об одном нелинейном уравнении с весовым оператором типа потенциала

\inbook Труды восьмой Всероссийской научной конференции с международным участием. \rm Часть~3

\bookinfo Дифференциальные уравнения и краевые задачи

\serial Матем. моделирование и краев. задачи

\yr 2011

\pages 24--27

\publ СамГТУ

\publaddr Самара

\misctransl 

\by Askhabov~S.\,N.

\paper On a nonlinear equation with a   potential-type weighing operator

\inbook Proceedings of the Eighth  All-Russian Scientific Conference with international participation. \rm Part~3

\serial Matem. Mod. Kraev. Zadachi

\yr 2011

\pages 24--27

\publ SamGTU

\publaddr Samara





\RBibitem{askhabov:4} 

\by Асхабов~С.\,Н.

\paper Приближенное решение нелинейных уравнений с~весовыми операторами типа потенциала

\jour Уфимск. матем. журн.

\yr 2011

\vol 3

\issue 4

\pages 8--13

\misctransl 

\by Askhabov~S.\,N.

\paper Approximate solution of nonlinear equations with weighted potential type operators

\jour Ufimsk. Mat. Zh.

\yr 2011

\vol 3

\issue 4

\pages 8--13



\end{thebibliography}





\noindent

\hskip6.5cm \thedate \par\noindent

\hskip6.5cm\therevdate



\vspace{-8mm}





\makeengtitle



\newpage



%\end{document}

